\chapter{Rôles et permissions}
\label{ch:roles}

mybibli a trois rôles\index{rôles} utilisateur, ordonnés par
privilège :

\begin{enumerate}
  \item \textbf{Anonyme}\index{anonyme} — quiconque arrive sur le
        site sans cookie de session. Accès en lecture seule au
        catalogue. Ne peut ni prêter, ni éditer, ni supprimer.
  \item \textbf{Bibliothécaire}\index{bibliothécaire} — utilisateur
        authentifié de rôle \texttt{librarian}. Catalogage au
        quotidien, prêts, retours, déplacements de lieu. Ne peut
        ni gérer les utilisateurs ni changer les paramètres
        système.
  \item \textbf{Administrateur}\index{administrateur} — utilisateur
        authentifié de rôle \texttt{admin}. Tout ce qu'un
        bibliothécaire peut faire, plus le panneau \texttt{/admin}
        : utilisateurs, données de référence, corbeille,
        paramètres système.
\end{enumerate}

\section{Qui peut faire quoi}

\begin{longtable}{@{} l c c c @{}}
\toprule
\textbf{Action} & \textbf{Anon} & \textbf{Bibl.} & \textbf{Admin} \\
\midrule
\endhead
Parcourir / rechercher dans le catalogue & oui & oui & oui \\
Voir les pages de titre et de volume     & oui & oui & oui \\
Voir l'arbre des lieux                   & oui & oui & oui \\
\midrule
Scanner / cataloguer un nouveau titre    & non & oui & oui \\
Ajouter ou retirer un volume             & non & oui & oui \\
Déplacer un volume entre lieux           & non & oui & oui \\
Éditer les métadonnées d'un titre        & non & oui & oui \\
Créer / éditer les emprunteurs           & non & oui & oui \\
Enregistrer / retourner des prêts        & non & oui & oui \\
Éditer l'arbre des lieux                 & non & oui & oui \\
\midrule
Ouvrir \texttt{/admin}                   & non & non & oui \\
Créer / désactiver des utilisateurs      & non & non & oui \\
Éditer les données de référence          & non & non & oui \\
Restaurer depuis la corbeille            & non & non & oui \\
Supprimer définitivement de la corbeille & non & non & oui \\
Éditer les paramètres système            & non & non & oui \\
\bottomrule
\end{longtable}

\section{Sessions et timeout d'inactivité}

Après la connexion, mybibli pose un cookie de
session\index{cookie de session} qui voyage avec chaque requête
suivante. Le cookie est :

\begin{itemize}
  \item \texttt{HttpOnly} — non exposé à JavaScript ;
  \item \texttt{SameSite=Lax} — envoyé sur les navigations
        same-site de haut niveau mais pas sur les requêtes
        cross-site ;
  \item \texttt{Secure} (uniquement si
        \texttt{MYBIBLI\_COOKIE\_SECURE} est positionné à
        \texttt{true} — voir chapitre~\ref{ch:installation}).
\end{itemize}

La session a un \textbf{timeout d'inactivité}\index{timeout
d'inactivité} (défaut 30 minutes, configurable depuis
\texttt{/admin > System > Loans}). Chaque requête authentifiée
réinitialise le compte à rebours ; si aucune requête n'arrive dans
la fenêtre, la session est effacée côté serveur et la requête
suivante redirige vers la page de connexion.

Un petit \emph{toast de keep-alive} apparaît dans un coin de l'écran
environ deux minutes avant l'expiration — cliquez-le pour
prolonger la session sans quitter la page courante.

\section{Garde du dernier admin actif}

mybibli refuse de laisser l'installation sans aucun admin actif.
La garde s'applique à deux actions :

\begin{itemize}
  \item \textbf{Désactiver un utilisateur.} Si l'utilisateur est le
        seul admin actif, la requête renvoie un 409 conflict et la
        désactivation est refusée.
  \item \textbf{Supprimer définitivement un utilisateur depuis la
        corbeille.} Même règle.
\end{itemize}

L'auto-désactivation est également bloquée sans condition (vous ne
pouvez pas désactiver le compte avec lequel vous êtes connecté).

\begin{warning}
La garde ne compte que les admins \emph{actifs}. Si vous
désactivez en masse des admins et que la garde se déclenche,
défaites en réactivant quelqu'un depuis l'onglet corbeille. Il n'y
a aucun chemin de récupération via l'assistant ou la base sans
SQL manuel.
\end{warning}

\section{Réinitialisation de mot de passe}

mybibli n'implémente pas l'envoi de courriels de réinitialisation
de mot de passe — les déploiements familiaux / petite communauté
tirent rarement profit de la plomberie email. Un admin peut
réinitialiser le mot de passe d'un autre utilisateur en :

\begin{enumerate}
  \item Ouvrant \texttt{/admin > Users}.
  \item Choisissant l'utilisateur, cliquant \textbf{Désactiver}.
  \item Cliquant \textbf{Réactiver} depuis l'onglet corbeille.
  \item Re-créant l'utilisateur avec un mot de passe frais.
\end{enumerate}

Un utilisateur qui connaît son mot de passe actuel peut le changer
depuis la page \texttt{/account}.

\section{Sessions anonymes}

Les visiteurs anonymes obtiennent aussi une ligne de session
(utilisée pour la protection CSRF et la livraison des mises à jour
HTMX en cours). Les sessions anonymes expirent après sept jours
d'inactivité et sont collectées par une tâche d'arrière-plan
quotidienne — elles n'affectent pas les utilisateurs connectés.
